% =============================================================
%  Chapter 1 — Introduction
%  Toward Continual Hamiltonian Learning in Unknown Environments
% =============================================================

\chapter{Introduction \\Toward Continual Hamiltonian Learning in Unknown Environments}
\label{chap:intro}

% ------------------------------------------------------------------
\section{The Core Challenge}
\label{sec:intro:challenge}
% ------------------------------------------------------------------

Autonomous agents operating in real environments must make decisions under
incomplete knowledge, changing local conditions, and evolving notions of risk.
A mobile robot sent into an unmapped building, an off-road vehicle traversing
terrain whose material properties are hidden beneath surface appearance, or a
deformable manipulator squeezing through a corridor whose geometry is revealed
only incrementally each of these settings shares a fundamental difficulty.
Successful behavior cannot be reduced to executing a fixed policy, following a
fixed planner, or optimizing a fixed reward signal.
Instead, the agent must \emph{continually revise how it moves through its state
space} as new structure is revealed.

The challenge is not merely to choose good actions.
It is to maintain a \textbf{decision process that can be corrected and enriched
as the environment changes}: geometrically, materially, and semantically.
Current paradigms handle parts of this challenge well, but none provides all
four properties simultaneously:
\begin{enumerate}
  \item \emph{structured dynamics} - a decision policy whose topology can be
        analysed and shaped rather than treated as a black box;
  \item \emph{continual adaptation} - the ability to revise behavior online as
        new information arrives, without forgetting previously learned structure;
  \item \emph{passivity-monitored updates} - correction mechanisms whose
        energy dissipation is tracked at runtime via a passivity surrogate,
        providing a checkable signal (not a formal certificate) that
        adaptation has not destabilized the field; and
  \item \emph{risk-sensitive objectives} - explicit reasoning about rare but
        catastrophic outcomes caused by latent material or hazard structure, not
        only about expected performance.
\end{enumerate}
We study how navigation in unknown environments can be cast as a
\textbf{continual Hamiltonian learning problem}—a principled framework that
jointly addresses all four properties by representing the agent's
decision-making as a \emph{structured phase-space dynamics} whose energy
landscape is progressively reshaped by geometry, constraints, and
semantic risk.

% ------------------------------------------------------------------
\section{What Existing Paradigms Hold Fixed}
\label{sec:intro:gap}
% ------------------------------------------------------------------

Understanding \emph{why} our approach is needed requires identifying a shared
structural gap across existing paradigms.
The claim is not that imitation learning, reinforcement learning, or classical
planning are insufficient in every sense each provides important and
well-studied properties.
The claim is more specific: each paradigm tends to \emph{hold fixed} something
that must remain adaptable in truly unknown environments.

\begin{figure}[t]
    \centering
    \includegraphics[width=0.95\textwidth]{figs/ch1_method_comparison.png}
    \caption{
    The first three columns illustrate common decision-making paradigms in robotics and machine learning:
    imitation learning uses a flat policy fitted to a fixed demonstration target;
    actor--critic reinforcement learning updates a policy through an external value evaluator under a fixed reward structure;
    and replanning methods repeatedly solve trajectories against an external cost map or planner template.
    In all three cases, evaluation is external to the evolving trajectory.
    By contrast, our framework represents policy as a phase-space flow
    \((\Pos_t,\Mom_t)\mapsto(\Pos_{t+1},\Mom_{t+1})\),
    where the cotangent variable \(\Mom_t\) acts as an internal evaluator, predictor, and guide.
    New information enters directly through \(\Mom_t\),
    which in turn shapes the force balance driving the next configuration update.
    This internal phase-space evaluator supports continual correction across segment, episode, and curriculum timescales, while preserving an explicit dissipative structure for stability.
    % The figure highlights the core thesis distinction: rather than relying on a fixed or external evaluator, continual Hamiltonian learning embeds sensing, evaluation, correction, and control inside the evolving dynamics.
    }
    \label{fig:ch1:method_comparison}
\end{figure}

\paragraph{Imitation learning.}
Methods such as behavioral cloning~\citep{pomerleau1991efficient} and the
DAgger family~\citep{ross2011reduction} achieve sample efficiency by fixing a
demonstration target and learning a reactive policy relative to that target.
Even interactive imitation methods~\citep{shi2025dagger} improve robustness to
distribution shift, but the resulting controller remains a \emph{policy
mapping}: it does not expose the qualitative structure of the induced
closed-loop dynamics nor guarantee stable behavior when the environment
departs from the demonstration regime.

\paragraph{Reinforcement learning.}
Deep RL methods such as PPO~\citep{schulman2017proximalpolicyoptimizationalgorithms} and
SAC~\citep{haarnoja2018softactorcriticoffpolicymaximum} are powerful continuous-control optimizers, but
they fix the reward structure and optimize a black-box value function relative
to it.
Performance degrades sharply under distribution shift~\citep{wijmans2019dd},
and the learned policy encodes no explicit language of stability, dissipation,
or safety invariants.
Hierarchical RL~\citep{vezhnevets2017feudalnetworkshierarchicalreinforcement} adds temporal
structure but does not resolve the fundamental issue: the optimizer still has
no direct handle on the \emph{topology} of the decision field being learned.

\paragraph{Inverse reinforcement learning and reward learning.}
IRL and preference-based methods~\citep{ziebart2008maximum,christiano2017deep}
learn objectives from demonstrations or comparisons, partially relaxing the
fixed-reward assumption.
Yet they still rely on a separate downstream optimization stage with a
fixed controller class, and the induced closed-loop dynamics need not be
topology-correct stably attracted to the goal, robustly repelled from obstacles across changing environments.

\paragraph{Classical planners and model-based methods.}
Graph search~\citep{hart1968formal}, sampling-based
planners~\citep{karaman2011sampling}, and MPC-style
approaches~\citep{williams2017mppi} provide explicit structure and, in some
cases, formal guarantees.
However, they require pre-specified cost representations and do not
themselves provide a stable, differentiable, continually correctable mechanism
for reshaping the decision dynamics as latent material or semantic properties
change.

\paragraph{LLM-based planners.}
Foundation-model planners~\citep{ahn2022saycan,liang2022code} offer impressive
semantic flexibility and goal generalization, but they govern
\emph{high-level} planning and still rely on separate low-level controllers
for stable continuous-time dynamics.
They do not themselves provide passivity, safety invariants, or risk-sensitive
field shaping at the level of the underlying motion.

Table~\ref{tab:paradigm_gap} summarizes this comparison along the four
axes identified above.

\begin{table}[htbp]
\centering
\caption{%
\textbf{What existing paradigms hold fixed versus what this thesis learns.}
Each paradigm solves a valuable subproblem.
Our contribution is to unify (i)--(iv) through continual Hamiltonian
learning.
$\checkmark$ = genuinely addressed; $\triangle$ = partial/monitored; $\times$ = not addressed.
Column (iii) for our work reflects passivity-regularized and
passivity-monitored updates; formal update-stability theorems are identified
as a contribution in Chapter~\ref{chap:limitations}.
}
\label{tab:paradigm_gap}
\small
\setlength{\tabcolsep}{6pt}
\renewcommand{\arraystretch}{1.25}
\resizebox{\textwidth}{!}{
\begin{tabular}{lcccc}
\toprule
\textbf{Paradigm} &
\makecell{\textbf{(i)}\\\textbf{Structured}\\\textbf{dynamics}} &
\makecell{\textbf{(ii)}\\\textbf{Continual}\\\textbf{adaptation}} &
\makecell{\textbf{(iii)}\\\textbf{Passivity-}\\\textbf{monitored updates}} &
\makecell{\textbf{(iv)}\\\textbf{Risk-sensitive}\\\textbf{objectives}} \\
\midrule
Imitation learning \citep{ross2011reduction,shi2025dagger}
  & \none & \partially & \none & \none \\
Deep RL \citep{schulman2017proximalpolicyoptimizationalgorithms,haarnoja2018softactorcriticoffpolicymaximum}
  & \none & \partially & \none & \none \\
Constrained / Safe RL \citep{achiam2017constrained,chow2015risk}
  & \none & \partially & \partially & \full \\
Inverse RL \citep{ziebart2008maximum,christiano2017deep}
  & \none & \none & \none & \none \\
Classical planning / MPC \citep{hart1968formal,williams2017mppi}
  & \partially & \none & \partially & \none \\
Port-Hamiltonian networks \citep{desai2021porthamiltonian,roth2025stable}
  & \full & \none & \full & \none \\
Continual / meta-RL \citep{rolnick2019experience,finn2017maml}
  & \none & \full & \none & \none \\
\midrule
\textbf{This thesis (Continual Hamiltonian Learning)}
  & \full & \full & \partially & \full \\
\bottomrule
\end{tabular}}
\end{table}

% ------------------------------------------------------------------
\section{The Thesis Viewpoint: Continual Hamiltonian Learning}
\label{sec:intro:viewpoint}
% ------------------------------------------------------------------

We adopt the following viewpoint:
\begin{quote}
\itshape
Rather than learning a fixed policy, reward, or planner, we seek to learn
a \textbf{structured decision dynamics}, a vector field over phase space that
can be \textbf{continually reshaped} by new observations.
\end{quote}
The key idea is to represent the agent's policy not as a mapping from
observations to actions, but as the deterministic \emph{flow map} of a
Hamiltonian vector field over the cotangent bundle $T^*\mathcal{Q}$:
\begin{equation}
  z_{t+1} = \Phi^{\Delta t}_{H}(z_t),
  \qquad
  \pi_H(z_t) := \Pi_{\mathcal{Q}}\!\left(\Phi^{\Delta t}_H(z_t)\right),
  \label{eq:intro:policy_as_flow}
\end{equation}
where $z_t = (\vec{q}_t, p_t) \in T^*\mathcal{Q}$,
$\Pi_{\mathcal{Q}}$ projects to configuration, and
$H : T^*\mathcal{Q} \to \mathbb{R}$ is a structured energy function
whose additive factors encode goal attraction, obstacle repulsion, risk, and dissipation.
Adapting the policy means \emph{reshaping the energy landscape} by adjusting
coefficients, adding new force channels, updating the dissipative port rather
than relearning a flat action map from scratch.

This perspective yields a natural language for all four desiderata:
\begin{enumerate}
  \item \emph{Structured dynamics.}
        The vector field topology is explicit: goal-attracting basins,
        obstacle-repelling barriers, and dissipative channels are
        interpretable terms in $H$.
  \item \emph{Continual adaptation.}
        New observations update the coefficients of energy terms, injecting
        minimal correction rather than catastrophic relearning.
  \item \emph{Passivity-monitored updates.}
        Passivity regularization and the port-Hamiltonian energy balance provide
        checkable conditions and empirical monitors for whether an adaptation
        step degrades stability~\citep{desai2021porthamiltonian,roth2025stable}.
        Formal update-stability theorems (sufficient conditions for
        passivity preservation under coefficient change) are identified
        as a contribution target in Chapter~\ref{chap:limitations}.
  \item \emph{Risk-sensitive objectives.}
        Tail-risk terms (e.g., CVaR over material cost distributions)
        enter as additional energy factors with their own force channels and
        gradient pathways~\citep{chow2015risk}.
\end{enumerate}

% ------------------------------------------------------------------
\section{The Thesis Arc: A Hierarchy of Complexity}
\label{sec:intro:arc}
% ------------------------------------------------------------------

We develop the continual Hamiltonian learning viewpoint
through two major stages, each solving a harder version of the same
underlying problem by enriching the energy landscape along well-defined axes.
Figure~\ref{fig:intro:hierarchy} illustrates this hierarchy.

\begin{figure}[htbp]
\centering
\begin{tikzpicture}[
  font=\small\sffamily,
  % --- node styles ---
  stage/.style={
    rectangle, rounded corners=5pt,
    draw=#1!70!black, fill=#1!10,
    line width=1.2pt,
    minimum width=11.2cm, minimum height=1.5cm,
    align=left, inner sep=8pt,
    text=black
  },
  stageno/.style={
    circle, draw=#1!80!black, fill=#1!25,
    line width=1pt,
    minimum size=22pt,
    font=\bfseries\small\sffamily
  },
  axis/.style={
    rectangle, rounded corners=3pt,
    draw=gray!60, fill=gray!8,
    inner sep=4pt, align=center,
    font=\scriptsize\sffamily
  },
  arrow/.style={-{Stealth[length=6pt]}, line width=1.2pt, color=#1},
  dasharrow/.style={-{Stealth[length=5pt]}, dashed, line width=0.9pt, color=gray!60},
  addlabel/.style={font=\scriptsize\sffamily\itshape, text=gray!70!black},
]

% -------------------------------------------------------
%  Stage boxes (stacked vertically, left-aligned content)
% -------------------------------------------------------

% Stage 0 — foundation
\node[stage=teal] (s0) at (0,0) {%
  \hspace{26pt}%
  \begin{minipage}{9.8cm}
    \textbf{Stage 1 — Geometry-Only Foundation (GRL-SNAM)}\\[2pt]
    \textit{Energy:} $H = H_{\mathrm{goal}} + H_{\mathrm{obs}} + H_{\mathrm{diss}}$\\
    \textit{Adaptation:} offline meta-training + online geometric alignment\\
    \textit{Limitation:} geometry-blind to latent material/risk structure
  \end{minipage}%
};
\node[stageno=teal] at ($(s0.west)+(14pt,0)$) {\textcolor{teal!80!black}{1}};

% Stage 1 — material-aware
\node[stage=orange, below=14pt of s0] (s1) {%
  \hspace{26pt}%
  \begin{minipage}{9.8cm}
    \textbf{Stage 2 — Material-Aware Risk-Sensitive Learning}\\[2pt]
    \textit{Energy:} $H = H_{\mathrm{goal}} + H_{\mathrm{obs}} + H_{\mathrm{mat}}^{\mathrm{soft}}
                        + H_{\mathrm{haz}}^{\mathrm{hard}} + H_{\mathrm{diss}}$\\
    \textit{Adaptation:} segment/episode/curriculum gradient updates\\
    \textit{Adds:} new force channels, tail-risk (CVaR) objective, oracle material maps
  \end{minipage}%
};
\node[stageno=orange] at ($(s1.west)+(14pt,0)$) {\textcolor{orange!80!black}{2}};

% Stage 2 — synthesis
\node[stage=violet, below=14pt of s1] (s2) {%
  \hspace{26pt}%
  \begin{minipage}{9.8cm}
    \textbf{Synthesis — Continual Hamiltonian Learning Framework}\\[2pt]
    \textit{Unifies:} stages 1--2 as a hierarchy of energy enrichment\\
    \textit{Shows:} each new term $\Rightarrow$ new force channel, coefficient,
                    loss, update law\\
    \textit{Vision:} self-improving learner in unknown environments
  \end{minipage}%
};
\node[stageno=violet] at ($(s2.west)+(14pt,0)$) {\textcolor{violet!80!black}{$\Sigma$}};

% -------------------------------------------------------
%  Right-side axis labels
% -------------------------------------------------------
\node[axis, right=18pt of s0] (a0)
  {\textbf{Geometry}\\\textbf{sensing}};
\node[axis, right=18pt of s1] (a1)
  {\textbf{Material}\\\textbf{risk}};
\node[axis, right=18pt of s2] (a2)
  {\textbf{Continual}\\\textbf{learning}};

% -------------------------------------------------------
%  Vertical progression arrow
% -------------------------------------------------------
\draw[arrow=teal!60!black]
  ($(s0.south west)!0.5!(s0.south east)+(0,-2pt)$) --
  ($(s1.north west)!0.5!(s1.north east)+(0,2pt)$)
  node[midway, left=4pt, addlabel] {enriched energy $+$ new objectives};
\draw[arrow=orange!70!black]
  ($(s1.south west)!0.5!(s1.south east)+(0,-2pt)$) --
  ($(s2.north west)!0.5!(s2.north east)+(0,2pt)$)
  node[midway, left=4pt, addlabel] {unified hierarchy $+$ synthesis};

% -------------------------------------------------------
%  Horizontal axis arrows (right side)
% -------------------------------------------------------
\draw[arrow=teal!50!black]   (s0.east) -- (a0.west);
\draw[arrow=orange!60!black] (s1.east) -- (a1.west);
\draw[arrow=violet!60!black] (s2.east) -- (a2.west);

% -------------------------------------------------------
%  "What is enriched" annotation box at top
% -------------------------------------------------------
\node[
  rectangle, rounded corners=3pt,
  draw=gray!40, fill=white,
  inner sep=5pt,
  font=\scriptsize\sffamily,
  text width=6.4cm,
  above=16pt of s0,
  align=center
] (legend) {%
  \textbf{At each stage, adding one energy term induces:}\\[2pt]
  new force channel $\;\cdot\;$ new coefficient\\
  new objective term $\;\cdot\;$ new gradient pathway\\
  new update law across timescales
};
\draw[dasharrow] (legend.south) -- ($(s0.north)+(0,2pt)$);

\end{tikzpicture}
\caption{%
  \textbf{The enrichment hierarchy of complexity.}
  Each stage enriches the phase-space energy landscape along a well-defined
  axis.
  Stage 1 establishes the
  geometry-only foundation with offline meta-training and online geometric
  alignment.
  Stage 2 extends this to material-aware,
  risk-sensitive continual learning by adding new energy terms, a CVaR
  tail-risk objective, and multi-timescale gradient updates.
  The synthesis unifies both stages as
  instances of a single enrichment principle and frames the broader
  continual Hamiltonian learning vision.
}
\label{fig:intro:hierarchy}
\end{figure}

\paragraph{Stage 1: geometry-only structured navigation (GRL-SNAM~\citep{ellendula2025grlsnam}).}
The first stage establishes the \emph{base case}: learning structured local
navigation under geometry-only sensing.
We formulate navigation as learning a reduced Hamiltonian
$H : T^*\mathcal{Q} \to \mathbb{R}$ whose energy factors capture goal
attraction, obstacle repulsion, and dissipation.
The policy is the time-$\Delta t$ flow map $\Phi^{\Delta t}_H$, learned through a three-module
architecture operating at different timescales sensing, frame-level motion,
and shape control.
Offline meta-training identifies the family of admissible control objectives;
online geometric alignment updates energy coefficients as new obstacles are
revealed.
GRL-SNAM demonstrates that geometry-aware energy shaping significantly
outperforms standard RL baselines
and classical planners on deformable robot navigation in unknown environments.

At the same time, Stage~1 reveals clear limitations that motivate the next
stage: the energy landscape is \emph{geometry-only}.
It cannot reason about latent material properties like terrain compliance, hazard
density, absorption coefficients that affect traversal cost and risk in
ways invisible to geometric sensing alone.

\paragraph{Stage 2: material-aware risk-sensitive continual Hamiltonian learning.}
The second stage enriches the energy landscape with two new factors:
a \emph{soft material-risk term} $H_{\mathrm{mat}}^{\mathrm{soft}}$ encoding
expected material cost, and a \emph{hard hazard barrier}
$H_{\mathrm{haz}}^{\mathrm{hard}}$ that acts as a steep repulsive potential
over high-risk regions, inspired by control barrier function
ideas~\citep{ames2016control} but without a formal forward-invariance certificate.
These additions induce new force channels, new coefficients to be learned,
new objective terms (including a CVaR tail-risk
penalty~\citep{chow2015risk}), and new update laws operating across segment,
episode, and curriculum timescales.
The result is a richer learner that adapts not only to geometric change but
to latent material structure revealed through experience.

\paragraph{Synthesis.}
We conclude by synthesizing both stages as instances of a single
enrichment principle: each new energy term induces a new force channel, a
new coefficient subspace, a new objective term, a new gradient pathway, and
a new update law.
This hierarchy is the backbone of a broader continual Hamiltonian learning
vision—a framework for agents that progressively enrich their decision
dynamics as the unknown environment reveals its structure.

% ------------------------------------------------------------------
\section{Why the Hamiltonian Viewpoint?}
\label{sec:intro:why}
% ------------------------------------------------------------------

The choice of Hamiltonian structure is not aesthetic.
It provides five concrete advantages for the problem studied here.

\paragraph{1. Interpretable topology.}
The energy factorization makes the qualitative topology of the decision
field explicit: basins of attraction around goals, repulsive walls around
obstacles, and dissipative channels that remove excess kinetic energy.
This topology can be verified and corrected independently of the policy
weights~\citep{greydanus2019hamiltonian}.

\paragraph{2. Stability language.}
Port-Hamiltonian systems admit a natural stability analysis through passivity
and energy dissipation~\citep{desai2021porthamiltonian}.
When a new energy term is added, passivity monitoring provides a checkable
signal—whether the adaptation step increased or decreased total system
energy—something that flat RL policies do not naturally
expose~\citep{roth2025stable}.
This is not a formal Lyapunov certificate but a practical indicator
that guides the curriculum (§\ref{sec:ch4:loops}).

\paragraph{3. Compositionality.}
Complex behavior emerges from composing interpretable energy terms that can
be reweighted or extended without relearning from scratch.
This is precisely the mechanism needed for continual adaptation: the agent
adds terms to its existing energy landscape rather than discarding
learned structure~\citep{rolnick2019experience}.

\paragraph{4. Differentiable geometry.}
The symplectic structure of $T^*\mathcal{Q}$ preserves geometric invariants
(energy, volume, momentum) through numerical
integration~\citep{reich1996symplectic}, reducing long-horizon drift that
plagues black-box policy networks on continuous-control tasks.

\paragraph{5. Natural coupling to sensing.}
New sensor observations—geometric clearance, material spectral response,
traversability estimates~\citep{rugd2019,rellis2020}—enter as updates to
energy coefficients, directly reshaping the force field without requiring
a separate reward-engineering step.

% ------------------------------------------------------------------
\section{Scope and What This Thesis Does Not Claim}
\label{sec:intro:scope}
% ------------------------------------------------------------------

We make focused claims and do not attempt to replace all existing
approaches.
Three boundaries are important to state explicitly.

\emph{Not a general continual learning algorithm.}
The framework is developed and validated in the context of robot navigation
in unknown 2D/3D environments.
Generalizing to arbitrary task structures, object manipulation, or social
interaction is future work.

\emph{Not a learned perception front-end.}
We assume that a sensor pipeline provides geometry and
material-context observations.
Learning the perception co-jointly with the energy landscape is an
important extension~\citep{rosinol2020kimera} but is not the primary focus.

\emph{Principled framework, not universal benchmark winner.}
The claim is that continual Hamiltonian learning is a \emph{principled
alternative} to black-box policy learning in unknown environments: one that
provides structure, stability, and compositional enrichment that existing
methods lack.
The experimental chapters validate this claim on the problem settings studied;
they do not claim superiority across all benchmarks.

% ------------------------------------------------------------------
\section{Thesis Contributions}
\label{sec:intro:contributions}
% ------------------------------------------------------------------

The main contributions of this work are:

\begin{enumerate}
  \item \textbf{GRL-SNAM: geometry-only Hamiltonian navigation}
        A geometric RL framework that casts simultaneous navigation and
        mapping as stagewise Hamiltonian optimization, with a three-module
        architecture, offline meta-training, and online geometric alignment.
        \textit{\textbf{Accepted at ICLR 2026.}}

  \item \textbf{Material-aware risk-sensitive continual Hamiltonian learning}
        An extension that enriches the energy landscape with soft material-risk
        and hard hazard-barrier terms, introduces CVaR tail-risk objectives,
        and defines multi-timescale gradient update laws across segment,
        episode, and curriculum loops.

  \item \textbf{A hierarchy-of-complexity analysis}
        A unified view showing that each stage of enrichment follows the
        same structural pattern: new energy term $\Rightarrow$ new force
        channel, coefficient, objective, and update law and framing both
        stages as instances of a single continual Hamiltonian learning
        principle.

  \item \textbf{Empirical validation and failure analysis}
        Comparative experiments on deformable robot navigation and
        material-rich terrain, analyzing both quantitative performance and
        qualitative failure modes as a function of energy complexity.
\end{enumerate}

% ------------------------------------------------------------------
\section{Thesis Roadmap}
\label{sec:intro:roadmap}
% ------------------------------------------------------------------
The remainder of this document is organized as follows.

\textbf{Chapter~2 (Background)} unifies the mathematical tools shared
across both stages: optimal control and the Pontryagin viewpoint,
port-Hamiltonian dynamics, symplectic integration, barrier methods,
risk-sensitive objectives (CVaR / CMDPs), and
continual/curriculum learning.

\textbf{Chapter~3 (GRL-SNAM)} presents the geometry-only foundation.
It introduces the three-module Hamiltonian architecture, the offline
meta-training procedure, and the online geometric alignment mechanism.
The chapter ends by identifying the limitations that motivate Stage~2.

\textbf{Chapter~4 (Material-Aware Continual Hamiltonian Learning)}
introduces the enriched energy landscape with material-risk and
hazard-barrier terms, the CVaR tail-risk objective, and the
multi-timescale gradient update laws.
This chapter constitutes the primary new contribution of this work.

\textbf{Chapter~5 (Hierarchy of Complexity)} synthesizes Chapters~3
and~4 by formalizing the enrichment principle and showing how each
added energy term induces new structure in the force field, objective,
and update law.

\textbf{Chapter~6 (Experimental Evaluation)} presents comparative
results across both stages, analysing performance, failure modes, and
the qualitative impact of each energy enrichment step.

\textbf{Chapters~7--8 (Limitations, Future Work, and Conclusion)}
discuss open problems including perception co-training, stronger
passivity guarantees, and richer coverage theory and synthesize the
enrichment arc.